% ==================================================
% components_base.tex
% 通用组件 / 通用排版规则（基础层, layered params）
% Parameter access: \csname HB<key>\endcsname
% ==================================================

\usepackage{array}
\usepackage{colortbl}
\usepackage{graphicx}
\usepackage{needspace}

% ------------------------------
% Raw LaTeX component helpers
% ------------------------------
\makeatletter

\providecommand{\HBIfBlankTF}[3]{%
  \edef\HB@blanktest{\detokenize{#1}}%
  \ifx\HB@blanktest\@empty
    #2%
  \else
    #3%
  \fi
}

\providecommand{\HBImagePlaceholder}[2]{%
  \makebox[#1][c]{\rule{0pt}{#2}}%
}

\providecommand{\HBImageOrPlaceholder}[3]{%
  \HBIfBlankTF{#3}{%
    \HBImagePlaceholder{#1}{#2}%
  }{%
    \IfFileExists{#3}{%
      \includegraphics[width=#1,height=#2,keepaspectratio]{#3}%
    }{%
      \HBImagePlaceholder{#1}{#2}%
    }%
  }%
}

\makeatother

% ------------------------------
% List base (enumitem)
% ------------------------------
\setlist[itemize]{
  leftmargin=\csname HBcomp_list_leftmargin\endcsname,
  labelsep=\csname HBcomp_list_labelsep\endcsname,
  itemsep=\csname HBcomp_list_itemsep\endcsname,
  topsep=\csname HBcomp_list_topsep\endcsname,
  parsep=\csname HBcomp_list_parsep\endcsname,
  partopsep=\csname HBcomp_list_partopsep\endcsname,
  label=\raisebox{\csname HBcomp_list_bullet_raise\endcsname}{\csname HBcomp_list_bullet_symbol\endcsname}
}

% Nested manual items are a distinct dash-led level in the production
% master.  Keep the hierarchy semantic so Sphinx, LaTeX, and IDML can all
% preserve it instead of flattening child items into parent prose.
\setlist[itemize,2]{
  leftmargin=\csname HBcomp_sublist_leftmargin\endcsname,
  labelsep=\csname HBcomp_sublist_labelsep\endcsname,
  itemsep=\csname HBcomp_sublist_itemsep\endcsname,
  topsep=\csname HBcomp_sublist_topsep\endcsname,
  parsep=0pt,
  partopsep=0pt,
  label=\csname HBcomp_sublist_bullet_symbol\endcsname
}

% ------------------------------
% Table base
% ------------------------------
\arrayrulewidth=\csname HBcomp_table_inner_rule\endcsname
\setlength{\tabcolsep}{\csname HBcomp_table_tabcolsep\endcsname}

\newtcolorbox{tableframe}{
  enhanced,
  colback=white,
  colframe=LineK40,
  arc=\csname HBcomp_table_outer_arc\endcsname,
  boxrule=\csname HBcomp_table_outer_rule\endcsname,
  left=0mm,right=0mm,top=0mm,bottom=0mm,
}

% ------------------------------
% Shared spec-like table style
% - Used by spec page and symbols page
% ------------------------------
\expandafter\providecommand\csname HBcomp_spec_table_left_ratio\endcsname{0.33}
\expandafter\providecommand\csname HBcomp_spec_table_tabcolsep\endcsname{3.2pt}
\expandafter\providecommand\csname HBcomp_spec_table_row_stretch\endcsname{1.10}

% Left-column tint is parameterized; the V2.0 JE-1000F master
% (2026-06-05) draws spec/LCD tables WITHOUT a tinted left column, so
% the default is 0 (white). Set comp_spec_table_left_tint in
% layout_params.csv to restore a tint for older-style targets.
\expandafter\providecommand\csname HBcomp_spec_table_left_tint\endcsname{0}
\colorlet{HBSharedTableLeftBg}{black!\csname HBcomp_spec_table_left_tint\endcsname}
% Table header row tint (Symbol/Meaning header bars in the V2.0 master)
\colorlet{HBTableHeadBg}{black!8}
\colorlet{HBSharedTableFrameColor}{black}

\providecommand{\HBsharedXcolumn}[1]{m{#1}}
\providecommand{\HBsharedTableLeftRatio}{0.33}
\providecommand{\HBsharedTableTabcolsep}{3.2pt}
\expandafter\providecommand\csname HBcomp_table_text_indent\endcsname{5.2pt}

\providecommand{\HBSetSharedTableGeometry}[2]{%
  \def\HBsharedTableLeftRatio{#1}%
  \def\HBsharedTableTabcolsep{#2}%
}

\newtcolorbox{hbsharedtableframe}{
  enhanced,
  clip upper,
  colback=white,
  colframe=HBSharedTableFrameColor,
  arc=\csname HBcomp_table_outer_arc\endcsname,
  boxrule=\csname HBcomp_table_outer_rule\endcsname,
  boxsep=0mm,
  left=0mm,right=0mm,top=0mm,bottom=0mm,
  before skip=0mm,after skip=0mm,
  overlay unbroken={
    \begin{tcbclipinterior}
      \fill[HBSharedTableLeftBg]
        (interior.south west)
        rectangle
        ([xshift=\dimexpr\HBsharedTableLeftRatio\linewidth+\HBsharedTableTabcolsep\relax]interior.north west);
    \end{tcbclipinterior}
  },
}

\newenvironment{HBSharedDataTable}[2]{%
  \HBSetSharedTableGeometry{#1}{#2}%
  \begin{hbsharedtableframe}
  \begingroup
  \parindent=0pt\relax
  \arrayrulecolor{HBSharedTableFrameColor}
  \let\tabularxcolumn\HBsharedXcolumn
  \setlength{\tabcolsep}{#2}
  \renewcommand{\arraystretch}{\csname HBcomp_spec_table_row_stretch\endcsname}
  % Inner lines follow the spec rule: same color as frame, thinner than outer box.
  \arrayrulewidth=\dimexpr\csname HBcomp_table_outer_rule\endcsname/2\relax
  \ifdim\arrayrulewidth<0.24pt\arrayrulewidth=0.24pt\fi
}{%
  \endgroup
  \end{hbsharedtableframe}
}

% ------------------------------
% Note box
% ------------------------------
\newtcolorbox{notebox}{
  enhanced,
  colback=BgK05,
  colframe=BgK05,
  arc=\csname HBcomp_note_arc\endcsname,
  boxrule=0pt,
  left=\csname HBcomp_note_pad_lr\endcsname,
  right=\csname HBcomp_note_pad_lr\endcsname,
  top=\csname HBcomp_note_pad_tb\endcsname,
  bottom=\csname HBcomp_note_pad_tb\endcsname,
}

% ------------------------------
% Notice / note / tip / caution blocks for raw LaTeX renderers
% Signature: \HBNoticeBlock[kind]{label}{primary text}{secondary text}
% Supported kind values: notice, note, tip, caution.
% ------------------------------
\expandafter\providecommand\csname HBcomp_notice_before\endcsname{0.8mm}
\expandafter\providecommand\csname HBcomp_notice_after\endcsname{0.8mm}
\expandafter\providecommand\csname HBcomp_notice_pad_lr\endcsname{1.4mm}
\expandafter\providecommand\csname HBcomp_notice_pad_tb\endcsname{0.8mm}
\expandafter\providecommand\csname HBcomp_notice_arc\endcsname{0.9mm}
\expandafter\providecommand\csname HBcomp_notice_rule\endcsname{0.35pt}
\expandafter\providecommand\csname HBcomp_notice_label_width\endcsname{0.18\linewidth}
\expandafter\providecommand\csname HBcomp_notice_label_gap\endcsname{1.2mm}
\expandafter\providecommand\csname HBtype_notice_label_font_size\endcsname{6.8pt}
\expandafter\providecommand\csname HBtype_notice_label_font_leading\endcsname{7.4pt}
\expandafter\providecommand\csname HBtype_notice_body_font_size\endcsname{6.6pt}
\expandafter\providecommand\csname HBtype_notice_body_font_leading\endcsname{7.4pt}

\colorlet{HBNoticeBgnotice}{BgK05}
\colorlet{HBNoticeFramenotice}{LineK40}
\colorlet{HBNoticeLabelnotice}{TextK90}
\colorlet{HBNoticeBgnote}{BgK05}
\colorlet{HBNoticeFramenote}{LineK40}
\colorlet{HBNoticeLabelnote}{TextK90}
\colorlet{HBNoticeBgtip}{AccentBlue!7}
\colorlet{HBNoticeFrametip}{AccentBlue!60!black}
\colorlet{HBNoticeLabeltip}{AccentBlue!60!black}
\colorlet{HBNoticeBgcaution}{white}
\colorlet{HBNoticeFramecaution}{BrandDark}
\colorlet{HBNoticeLabelcaution}{BrandDark}

\makeatletter
\providecommand{\HBNoticeColor}[2]{%
  \ifcsname color@HBNotice#1#2\endcsname
    HBNotice#1#2%
  \else
    HBNotice#1note%
  \fi
}
\makeatother

\providecommand{\HBTypeNoticeLabel}[2]{%
  {\HBFontSemiBold
   \color{\HBNoticeColor{Label}{#1}}%
   \fontsize{\csname HBtype_notice_label_font_size\endcsname}
            {\csname HBtype_notice_label_font_leading\endcsname}\selectfont
   \RaggedRight #2}%
}

\providecommand{\HBTypeNoticeBody}[1]{%
  {\HBFontRegular
   \color{\csname HBtype_body_color\endcsname}%
   \fontsize{\csname HBtype_notice_body_font_size\endcsname}
            {\csname HBtype_notice_body_font_leading\endcsname}\selectfont
   \RaggedRight #1}%
}

\newcommand{\HBNoticeBlock}[4][notice]{%
  \par\noindent
  \begin{tcolorbox}[
    enhanced,
    breakable,
    colback=\HBNoticeColor{Bg}{#1},
    colframe=\HBNoticeColor{Frame}{#1},
    arc=\csname HBcomp_notice_arc\endcsname,
    boxrule=\csname HBcomp_notice_rule\endcsname,
    boxsep=0mm,
    left=\csname HBcomp_notice_pad_lr\endcsname,
    right=\csname HBcomp_notice_pad_lr\endcsname,
    top=\csname HBcomp_notice_pad_tb\endcsname,
    bottom=\csname HBcomp_notice_pad_tb\endcsname,
    before skip=\csname HBcomp_notice_before\endcsname,
    after skip=\csname HBcomp_notice_after\endcsname
  ]%
    \HBTypeNoticeLabel{#1}{#2}%
    \par\vspace{0.25mm}%
    \HBTypeNoticeBody{#3}%
    \HBIfBlankTF{#4}{}{%
      \par\vspace{0.35mm}%
      \HBTypeNoticeBody{#4}%
    }%
  \end{tcolorbox}%
}

% Four inline notice variants share one visual component. The source table
% supplies content; this component owns the rounded panel, split, and type.
\providecommand{\HBCalloutBlock}[4]{%
  \par\noindent
  \begin{tcolorbox}[
    enhanced,
    colback=BgK05,
    colframe=BgK05,
    arc=\csname HBcomp_tip_arc\endcsname,
    boxrule=\csname HBcomp_callout_rule\endcsname,
    boxsep=0mm,
    leftupper=0mm,rightupper=0mm,
    leftlower=0mm,
    rightlower=\csname HBcomp_tip_pad_lr\endcsname,
    top=\csname HBcomp_caution_pad_tb\endcsname,
    bottom=\csname HBcomp_caution_pad_tb\endcsname,
    height from={#4} to {\textheight},
    valign=center,
    valign lower=center,
    sidebyside,
    sidebyside align=center,
    sidebyside gap=\csname HBcomp_callout_body_inset\endcsname,
    lefthand width=#3,
    before skip=1.2mm,
    after skip=1.2mm,
    underlay={
      \fill[white]
        ([xshift=\csname HBcomp_callout_label_inset\endcsname,
          yshift=\csname HBcomp_callout_label_inset\endcsname]frame.south west)
        rectangle
        ([xshift=#3,
          yshift=-\csname HBcomp_callout_label_inset\endcsname]frame.north west);
      \node[anchor=base,inner sep=0pt] at
        ([xshift=\dimexpr(#3+\csname HBcomp_callout_label_inset\endcsname)/2\relax,
          yshift=\csname HBcomp_tip_label_baseline_shift\endcsname]frame.west)
        {\HBFontBold
         \fontsize{\csname HBtype_tip_label_font_size\endcsname}
                  {\csname HBtype_tip_label_font_leading\endcsname}\selectfont
         #1};
    }
  ]%
    \begin{minipage}[c]{\linewidth}%
      \mbox{}%
    \end{minipage}%
    \tcblower
    \HBFontMedium
    \addfontfeatures{FakeStretch=\csname HBtype_tip_body_horizontal_scale\endcsname}%
    \fontsize{\csname HBtype_tip_body_font_size\endcsname}
             {\csname HBtype_tip_body_font_leading\endcsname}\selectfont
    \setlength{\parskip}{0pt}%
    \RaggedRight #2\par
  \end{tcolorbox}%
  \par
}
\providecommand{\HBCalloutBullet}[1]{%
  \par\noindent
  \hangindent=\dimexpr\csname HBcomp_callout_bullet_indent\endcsname
    +\csname HBcomp_callout_bullet_width\endcsname\relax
  \hangafter=1
  \hspace*{\csname HBcomp_callout_bullet_indent\endcsname}%
  \makebox[\csname HBcomp_callout_bullet_width\endcsname][l]{%
    \fontsize{\csname HBtype_callout_bullet_font_size\endcsname}
             {\csname HBtype_callout_bullet_font_size\endcsname}\selectfont\textbullet}%
  #1\par
}
\providecommand{\HBWarningBlock}[2]{%
  \HBCalloutBlock{#1}{#2}{\csname HBcomp_caution_label_width\endcsname}{0pt}%
}
\providecommand{\HBCautionBlock}[2]{%
  \HBCalloutBlock{#1}{#2}{\csname HBcomp_caution_label_width\endcsname}{0pt}%
}
\providecommand{\HBNoteBlock}[2]{%
  \HBCalloutBlock{#1}{#2}{\csname HBcomp_tip_label_width\endcsname}{0pt}%
}
\providecommand{\HBTipBlock}[2]{%
  \HBCalloutBlock{#1}{#2}{\csname HBcomp_tip_label_width\endcsname}%
    {\csname HBcomp_tip_height\endcsname}%
}
